\chapter*{Antes de começarmos...}
\addcontentsline{toc}{chapter}{Antes de começarmos...}
\chaptermark{Antes de começarmos...}

\lettrine[lines=4, lhang=0.1, lraise=0, loversize=0.2, findent=0.1em]{\textcolor{corTema}{P}}{ARABÉNS!} Se você chegou até aqui, é porque construiu uma base sólida e real de desenvolvimento de aplicações Web. Você sabe como o protocolo HTTP funciona, entende o ciclo de vida de uma requisição, sabe criar Servlets, montar páginas JSP com EL e JSTL, organizar o código em camadas usando os padrões DAO e MVC, conectar a aplicação a um banco de dados relacional, lidar com relacionamentos complexos e ainda dar vida à interface com JavaScript. Não é pouca coisa.

Mas o mundo do desenvolvimento Web não para. E uma das mudanças mais significativas das últimas décadas foi a popularização das arquiteturas baseadas em \textit{Web Services} RESTful. Nesse modelo, o servidor deixa de ser responsável por renderizar páginas HTML e passa a expor apenas dados --- normalmente no formato JSON --- por meio de \textit{endpoints} bem definidos. Quem cuida da interface é o cliente: um \textit{browser} rodando uma aplicação JavaScript, um aplicativo móvel, outro serviço no servidor. Essa separação entre \textit{backend} e \textit{frontend} trouxe uma flexibilidade enorme e hoje é, sem exagero, o padrão dominante no desenvolvimento de novas aplicações Web.

Você já viu um gostinho disso no capítulo~\ref{cap:javaScript}, quando usamos a Fetch API para fazer requisições assíncronas ao servidor e processar as respostas em JSON. Agora vamos dar um passo além: aprenderemos a construir o próprio \textit{backend} REST do zero, usando Servlets e a API JSON-B que você já conhece, sem depender de nenhum \textit{framework} adicional. Isso é importante, pois entender como as coisas funcionam por baixo dos panos é o que vai te diferenciar quando você precisar depurar um problema ou aprender um novo \textit{framework} no futuro.

Depois disso, o desafio final: reimplementar um sistema completo usando uma tecnologia diferente do Java --- uma forma de exercitar justamente essa capacidade de transferir o conhecimento que você construiu para um novo contexto.

Preparado? Vamos lá!
